\indent ರೀತಿಯ ಗಣವು ಬಂದಿರುವುದಿಲ್ಲ. (6) ಪೂರ್ವಾರ್ಧ ಹಾಗೂ ಉತ್ತರಾರ್ಧದಲ್ಲಿ ಆರನೆಯ ಮಾತ್ರಾಗಣವು ಮಧ್ಯಗುರುವುಳ್ಳ ಗಣವಾಗಿರುವುದು; ಇದು “ಲಲ” ಈ ರೀತಿಯಾಗಿಯೂ ಇರಬಹುದು, (7) ಪೂರ್ವಾರ್ಧದ ಅಂತ್ಯದಲ್ಲಿ ಮತ್ತು ಉತ್ತರಾರ್ಧದ ಅಂತ್ಯದಲ್ಲಿ ಗುರು ಬರುವ “ರ್ರ” ಅಥವಾ “\guru{}\guru{}” ರೀತಿಯ ಮಾತ್ರಾಗಣವಿರುವುದು.\par

\vspace{4pt plus 4pt minus 1pt}\noindent{\kbvsecfont\fontsize{11pt}{14.5pt}\selectfont\color{kbvaccentdark} (೨)\hspace{0.45em}ಷಟ್ಪದಿ:}\hspace{0.5em}ಇವು ಆರು ಸಾಲಿನ ಪದ್ಯಗಳು. ಷಟ್ಪದಿ ಪದ್ಯಗಳಲ್ಲಿ ಶರ ಷಟ್ಪದಿ, ಕುಸುಮ ಷಟ್ಪದಿ, ಭೋಗ ಷಟ್ಪದಿ, ಭಾಮಿನೀ ಷಟ್ಪದಿ, ಪರಿವರ್ಧಿನೀ ಷಟ್ಪದಿ, ವಾರ್ಧಕ ಷಟ್ಪದಿ ಎಂಬುದಾಗಿ ಆರು ಪ್ರಕಾರಗಳಿವೆ, ಇವುಗಳೆಲ್ಲವೂ ವಿಶಿಷ್ಟ ಲಕ್ಷಣಗಳಿಂದ ಕೂಡಿರುವುವು.\par

\indent ಷಟ್ಪದಿಗಳ ಲಕ್ಷಣಗಳು\par

\begin{kbvrule}
\noindent\kbvbullet\hspace{0.5em}ಕೆಳಗೆ ಕಂಡ ಭಾಮಿನೀ ಷಟ್ಪದಿಯ ಪದ್ಯಕ್ಕೆ ಪ್ರಸ್ತಾರ ಹಾಕಿ ಮಾತ್ರಾಗಣಗಳನ್ನು ವಿಂಗಡಿಸಿರುವುದನ್ನು ಗಮನಿಸಿ :
\end{kbvrule}

\indent \laghu{}\laghu{} \laghu{}\laghu{} \laghu{}\laghu{} \laghu{} \laghu{}\laghu{} \guru{}\laghu{}\laghu{} ಮಡದಿ ನಿಜ ನಿಲಯವನು ಬಿಟ್ಟಡಿ ಅ\par

\indent \laghu{}\laghu{} \laghu{}\laghu{} \laghu{}\laghu{}\laghu{}\laghu{}\laghu{} \guru{} \laghu{}\laghu{} ಗಡಿಗೆ ಪರಗೃಹದೊಳಗೆ ಬಾಯನು\par

\indent \laghu{}\laghu{}\laghu{}\laghu{} \laghu{}\laghu{}\laghu{}\laghu{}\guru{}\laghu{}\guru{}\laghu{}\laghu{}\guru{}\laghu{}\laghu{}\laghu{} \laghu{}\laghu{}\guru{}\par

\begin{kbvverse}
\noindent ಬಡಿದು ಮನೆಮನೆವಾರ್ತೆಯೆನ್ನದೆ ಬೀದಿಗಲಹವನು ।
\end{kbvverse}

\indent \laghu{}\laghu{}\laghu{}\laghu{} \laghu{}\laghu{}\laghu{}\laghu{}\guru{}\laghu{}\guru{}\guru{} ಒಡರಿಚುವ ಪತಿಯೊಬ್ಬನುಂಬೆಂ\par

\indent \laghu{}\laghu{}\laghu{} \guru{}\laghu{}\laghu{} \guru{}\laghu{} \guru{}\laghu{}\laghu{} ಬೆಡಹಿ ಕಾಣದ ದಿಟ್ಟೆ ಹತ್ತನು\par

\indent \laghu{}\laghu{}\laghu{} \laghu{}\laghu{} \guru{}\laghu{} \laghu{}\laghu{}\guru{}\laghu{}\laghu{} \laghu{}\laghu{}\laghu{}\laghu{}\guru{} \guru{}\par

\begin{kbvverse}
\noindent ಹಡೆದೊಡೆಯು ವರ್ಜಿಸುವುದುತ್ತಮ ಪುರುಷರುಗಳೆಂದ ॥
\end{kbvverse}

\indent (ಕರ್ಣಾಟ ಭಾರತ ಕಥಾಮಂಜರಿ)\par
